\section{The Live Voice Frontier}
\sectionkicker{VOICE FRONTIER}
\begin{wideflow}
{\Large\bfseries Gemini 3.8 Live、声で動く対話の最前線}\par
\smallskip
{\color{SurveyMuted}9月15日の2モデルを、音声対話の位置づけとベンダー側の数値の範囲で描く。}\par
\end{wideflow}

9月15日、Googleはライブ対話モデルの新系列としてGemini 3.8 LiveとGemini 3.8 Live Extended Thinkingを発表した\cite{geminilive}。聞き取りと発話を一つのモデルで担う音声対話の系列で、Liveは規模と費用効率に、Extended Thinkingは複雑な仕事向けの高い知能と多段階の推論に振るとされる。会話の途中で97言語を自動に行き来し、背景で道具やAPIを動かしながら会話を続ける振る舞いがうたわれる。Extended Thinkingは考えながら話す（「確認します」といった早い合図や、複数段階の進捗の実況）とされる。開発面ではGemini Live APIとGoogle AI Studio、企業向けにはGemini Enterpriseでの限定プレビュー、利用者向けにはGeminiアプリやSearch Live、Workspace（購読者向けのDocs・Gmail・Keep Live）が示される。生成音声にはすべてSynthIDの透かしが入るとされる。

測定として示されるのは、Extended Thinkingが音声対話品質指数で82.6の首位、tau-Voiceで68.6\%、銀行業務版で35.1\%、Big Bench Audioで97.7\%、Liveが音声エージェントの人気投票で2位、ServiceNowのEVA-Benchで複雑な業務の正確さと会話品質の両立が先頭に寄るとされることである\cite{geminilive}。いずれもベンダー側・引用元側のまとめであり、採点手法やサンプリング、比較条件の詳細な検証までは本号では行わない。Live APIやAI Studioの文書は頁上で結ばれるが、意味の読み込みまではしておらず、提供面は発表資料の範囲で扱う。

\begin{claimboundary}[測定と提供の境界]
測定値はベンダー側・引用元側の報告であり、独立した再現はない。文書頁の結びつきは提供の裏づけには使わない。対応言語や電話の範囲は資料だけでは輪郭が定まらず、言い切らない。
\end{claimboundary}

\begin{communitynote}[X / Community observation --- context only]
日本語圏の技術発信者によるGemini 3.8系列の組み込み試行（CodexやClaudeのパイプラインに載せる観察など）は、独立した利用者の文脈として残す\cite{w38x-gemini-jp}。公式の発信\cite{w38x-gemini-off}と合わせ、関心の向きを示す背景であり、公式の能力事実には繰り上げず、仕様や測定値の根拠にはしない。
\end{communitynote}
